% ============================================================
% Section 44: 金刚石的爱因斯坦模型——均方位移与零点能
% ============================================================
\newpage
\section{固体理论：金刚石的爱因斯坦模型——均方位移与零点能}
\label{sec:diamond-einstein}

\begin{modelbox}[题目：金刚石的爱因斯坦模型数值计算]
已知金刚石的爱因斯坦温度 $\Theta_E=1320\,\mathrm{K}$，碳原子质量 $m=19.92\times10^{-27}\,\mathrm{kg}$，晶格常数 $a=3.57\,\text{\AA}$。
\begin{enumerate}
    \item 估算室温 ($300\,\mathrm{K}$) 时，碳原子的方均根位移与原子间距的比值；
    \item 计算绝对零度下金刚石的单位体积零点振动能。
\end{enumerate}
\end{modelbox}

\begin{tipbox}[考点快速分析]
\begin{itemize}
    \item \textbf{三维爱因斯坦模型}：$3N$ 个独立谐振子，频率均为 $\omega_E$
    \item \textbf{均方位移公式}：$\langle r^2\rangle = \dfrac{3\hbar^2}{2mk_B\Theta_E}\coth\dfrac{\Theta_E}{2T}$
    \item \textbf{零点能}：$T=0$ 时每个模式仍有 $\dfrac12\hbar\omega_E$
    \item \textbf{金刚石结构}：晶胞含 8 个原子，最近邻距 $d=\dfrac{\sqrt3}{4}a$
\end{itemize}
\end{tipbox}

\begin{derivationbox}[标准解题过程]
\textbf{Step 1: 建立三维爱因斯坦模型的均方位移公式}

\begin{insightbox}[一维量子谐振子的均方位移推导]
对一维谐振子 $H=\frac{p^2}{2m}+\frac12m\omega_E^2x^2$，能量本征值 $E_n=(n+\frac12)\hbar\omega_E$。由维里定理，平均势能等于总能量的一半：
\[
\frac12m\omega_E^2\langle x^2\rangle_n = \frac12E_n = \frac12\Bigl(n+\frac12\Bigr)\hbar\omega_E.
\]
故在第 $n$ 个本征态中：
\begin{equation}
\langle x^2\rangle_n = \frac{(n+\frac12)\hbar}{m\omega_E}.
\end{equation}

对热平衡分布求平均，利用平均声子数 $\langle n\rangle=[e^{\hbar\omega_E/k_BT}-1]^{-1}$：
\begin{equation}
\langle x^2\rangle = \frac{(\langle n\rangle+\frac12)\hbar}{m\omega_E} = \frac{\hbar}{2m\omega_E}\coth\frac{\hbar\omega_E}{2k_BT}.
\end{equation}
最后一步利用了恒等式 $\coth x = \frac{e^x+e^{-x}}{e^x-e^{-x}} = 1+\frac{2}{e^{2x}-1}$ 以及 $\langle n\rangle=1/(e^{2x}-1)$（令 $x=\hbar\omega_E/2k_BT$）。
\end{insightbox}

在爱因斯坦模型中，每个原子在 $x,y,z$ 三个方向独立振动，频率均为 $\omega_E$。三维总均方位移为三个方向之和：
\begin{equation}
\boxed{\langle r^2\rangle = 3\langle x^2\rangle = \frac{3\hbar}{2m\omega_E}\coth\frac{\hbar\omega_E}{2k_BT} = \frac{3\hbar^2}{2mk_B\Theta_E}\coth\frac{\Theta_E}{2T}.}
\label{eq:diamond-msd-full}
\end{equation}

利用 $\coth x = 1+\dfrac{2}{e^{2x}-1}$，可写成等价形式：
\begin{equation}
\langle r^2\rangle = \frac{3\hbar^2}{mk_B\Theta_E}\Bigl(\frac{1}{e^{\Theta_E/T}-1}+\frac12\Bigr).
\label{eq:diamond-msd-alt}
\end{equation}
其中第一项为\textbf{热激发贡献}，第二项为\textbf{零点振动}。

\textbf{Step 2: 室温 ($300\,\mathrm{K}$) 下的方均根位移}

对金刚石，$T=300\,\mathrm{K}\ll\Theta_E=1320\,\mathrm{K}$，故 $e^{\Theta_E/T}=e^{4.4}\approx81.45\gg1$。

\textit{说明}：本题原解答中只取了热激发项（略去零点振动的 $1/2$）。严格来说 $T\ll\Theta_E$ 时零点振动占主导，不应略去；但以下按题目/教材的处理方式展示。

取热激发项：
\begin{equation}
\langle r^2\rangle \approx \frac{3\hbar^2}{mk_B\Theta_E}\cdot\frac{1}{e^{\Theta_E/T}-1}.
\end{equation}
代入常数：
\begin{align}
    \frac{3\hbar^2}{mk_B\Theta_E} &= \frac{3\times(1.0546\times10^{-34})^2}{(19.92\times10^{-27})(1.3807\times10^{-23})(1320)} \nonumber\\
    &\approx 9.19\times10^{-23}\,\mathrm{m^2} = 9.19\times10^{-3}\,\text{\AA}^2.
\end{align}
热激发因子：
\begin{equation}
\frac{1}{e^{4.4}-1} \approx \frac{1}{80.45} \approx 0.01243.
\end{equation}
因此
\begin{equation}
\langle r^2\rangle \approx 9.19\times10^{-23}\times0.01243 \approx 1.14\times10^{-24}\,\mathrm{m^2} = 1.14\times10^{-4}\,\text{\AA}^2.
\end{equation}
方均根位移
\begin{equation}
\sqrt{\langle r^2\rangle} \approx \sqrt{1.14\times10^{-4}}\,\text{\AA} \approx 1.07\times10^{-2}\,\text{\AA} \approx 0.0107\,\text{\AA}.
\end{equation}

\begin{insightbox}[金刚石结构的几何——8 个原子与最近邻距]
\textbf{1. 为什么晶胞含 8 个原子？}

金刚石结构由\textbf{两套 FCC 晶格}沿体对角线方向套构而成，套构位移为 $(\frac14,\frac14,\frac14)a$。
\begin{itemize}
    \item 单套 FCC 晶胞的原子数：$8\times\frac18$（角）$+6\times\frac12$（面）$=4$
    \item 两套 FCC 套构：$4\times2=8$ 个原子
\end{itemize}

\textbf{2. 为什么最近邻距 $d=\sqrt3\,a/4$？}

取第一套 FCC 角上的原子（坐标 $(0,0,0)$），它到第二套 FCC 中最近原子的距离为：
\[
d = \sqrt{\Bigl(\frac{a}{4}\Bigr)^2+\Bigl(\frac{a}{4}\Bigr)^2+\Bigl(\frac{a}{4}\Bigr)^2} = \sqrt{\frac{3a^2}{16}} = \frac{\sqrt3}{4}a.
\]
金刚石中每个碳原子与周围 4 个碳原子形成正四面体配位，键角 $109.5^\circ$，键长正是 $d=\sqrt3\,a/4$。
\end{insightbox}

金刚石为金刚石结构，最近邻原子间距：
\begin{equation}
\boxed{d = \frac{\sqrt3}{4}a = \frac{\sqrt3}{4}\times3.57\,\text{\AA} \approx 1.546\,\text{\AA}.}
\end{equation}
所求比值
\begin{equation}
\boxed{\frac{\sqrt{\langle r^2\rangle}}{d} \approx \frac{0.0107}{1.546} \approx 6.9\times10^{-3} \approx 0.69\%.}
\end{equation}

\textbf{Step 3: 绝对零度下的单位体积零点振动能}

$T=0$ 时热激发为零，但每个谐振子仍有零点能 $\dfrac12\hbar\omega_E$。晶体共 $N$ 个原子、$3N$ 个振动模式，总零点能
\begin{equation}
E_{\text{zero}} = 3N\cdot\frac12\hbar\omega_E = \frac32 Nk_B\Theta_E.
\end{equation}

金刚石结构晶胞体积 $a^3$ 内含 8 个原子，原子数密度 $N/V=8/a^3$。单位体积零点能
\begin{equation}
\boxed{E_0 = \frac{E_{\text{zero}}}{V} = \frac32\cdot\frac{8}{a^3}\cdot k_B\Theta_E = \frac{12k_B\Theta_E}{a^3}.}
\label{eq:diamond-zero-energy}
\end{equation}

代入数值：
\begin{align}
    a^3 &= (3.57\times10^{-10})^3 \approx 4.55\times10^{-29}\,\mathrm{m^3}, \\
    E_0 &= \frac{12\times(1.3807\times10^{-23})\times1320}{4.55\times10^{-29}} \nonumber\\
    &\approx \frac{2.186\times10^{-19}}{4.55\times10^{-29}} \approx 4.8\times10^9\,\mathrm{J/m^3}.
\end{align}
换算为 $\mathrm{cal/m^3}$（$1\,\mathrm{cal}=4.184\,\mathrm{J}$）：
\begin{equation}
\boxed{E_0 \approx 4.8\times10^9\,\mathrm{J/m^3} \approx 1.15\times10^9\,\mathrm{cal/m^3}.}
\end{equation}
\end{derivationbox}

\begin{formulabox}[答案汇总]
\begin{center}
\begin{tabular}{cl}
\toprule
问 & 结果 \\
\midrule
(1) 方均根位移/原子间距 & $\displaystyle \approx 0.69\%$ \\[8pt]
(2) 单位体积零点能 & $\displaystyle \approx 4.8\times10^9\,\mathrm{J/m^3}\;(=1.15\times10^9\,\mathrm{cal/m^3})$ \\
\bottomrule
\end{tabular}
\end{center}
\end{formulabox}

\begin{insightbox}[物理图像与概念深化]
\textbf{1. 原题解答是否严谨？——零点振动不可略}

严格公式 \eqref{eq:diamond-msd-alt} 中，$T=300\,\mathrm{K}$、$\Theta_E=1320\,\mathrm{K}$ 时：
\begin{itemize}
    \item 热激发项：$\dfrac{1}{e^{4.4}-1}\approx0.012$
    \item 零点项：$\dfrac12=0.5$
\end{itemize}
零点项是热激发项的 $\sim40$ 倍！若按完整公式计算：
\begin{align}
    \langle r^2\rangle_{\text{full}} &= \frac{3\hbar^2}{mk_B\Theta_E}\Bigl(0.012+\frac12\Bigr) \approx 9.19\times10^{-23}\times0.512 \nonumber\\
    &\approx 4.7\times10^{-23}\,\mathrm{m^2} = 4.7\times10^{-3}\,\text{\AA}^2,\\
    \sqrt{\langle r^2\rangle_{\text{full}}} &\approx \sqrt{4.7\times10^{-3}}\,\text{\AA} \approx 0.069\,\text{\AA}.
\end{align}
所求比值
\begin{equation}
\frac{\sqrt{\langle r^2\rangle_{\text{full}}}}{d} \approx \frac{0.069}{1.546} \approx 4.5\%.
\end{equation}

\textit{结论}：
\begin{itemize}
    \item 若问\textbf{总量子涨落}（绝对均方位移，含零点振动），正确答案应为 $\bm{\sim4.5\%}$，原题解答的 $0.69\%$ \textbf{遗漏了主导的零点贡献}，不够严谨。
    \item 若问\textbf{热振动位移}（仅由温度引起的、相对于 $T=0$ 平衡位置的额外偏移），则 $0.69\%$ 是合理的。这对应于估算热膨胀、Lindemann 判据等情境。
    \item 在 X射线衍射的 Debye--Waller 因子中，均方位移必须包含零点贡献；在讨论``热振幅''时，教材有时会只计热激发部分，但这需要明确说明。
\end{itemize}

\textbf{建议}：考试中若题目只说``估算均方位移''而未限定``热''，应给出完整结果（含零点能），或在答案中明确指出所作近似。

\textbf{2. 为什么金刚石的爱因斯坦温度如此高？}

$\Theta_E=1320\,\mathrm{K}$ 远高于室温，说明：
\begin{itemize}
    \item 碳原子质量轻 ($m\sim20\times10^{-27}\,\mathrm{kg}$)，惯性小
    \item C--C 共价键极强，有效弹簧常数 $k$ 很大
    \item 由 $\omega=\sqrt{k/m}$，轻质量 + 硬弹簧 → 高频率 → 高爱因斯坦温度
\end{itemize}
因此金刚石在室温下仍``冻结''在量子基态附近，热激发很弱，这正是它超高硬度和热稳定性的微观来源。

\textbf{3. 零点能的宏观意义}

单位体积零点能 $E_0\sim4.8\times10^9\,\mathrm{J/m^3}$ 是一个巨大的数字。作为对比：
\begin{itemize}
    \item 水的汽化热 $\sim2.3\times10^9\,\mathrm{J/m^3}$
    \item 金刚石的零点能与水的汽化热同数量级
\end{itemize}
零点能是量子力学不可避免的效应，即使在绝对零度也无法``关掉''晶格振动。对于轻元素（H、C）晶体，零点能尤为显著，甚至在某些情况下（如氦）会导致量子熔化。

\textbf{4. 爱因斯坦模型的数值计算技巧}

遇到``给定 $\Theta_E$、$m$、$a$，求具体数值''的问题：
\begin{enumerate}
    \item 先写出解析公式（如式 \eqref{eq:diamond-msd-full} 或 \eqref{eq:diamond-zero-energy}）
    \item 把所有常数一次性代入计算器，避免中间步骤舍入误差
    \item 注意单位统一（建议全部转换为国际单位制：$\text{\AA}\to10^{-10}\,\mathrm{m}$，$\mathrm{eV}\to1.6\times10^{-19}\,\mathrm{J}$）
    \item 最后检验数量级是否合理（金刚石原子间距 $\sim1.5\,\text{\AA}$，均方位移不可能超过此值）
\end{enumerate}
\end{insightbox}

\begin{thinkbox}[考场速记：爱因斯坦模型数值题的五步法]
\begin{enumerate}
    \item \textbf{写通式}：先写出不含数值的普遍公式（如 $E_0=12k_B\Theta_E/a^3$）
    \item \textbf{查结构}：确认晶胞原子数（金刚石 8、FCC 4、BCC 2、SC 1）
    \item \textbf{定方向}：均方位移注意是求一维分量还是三维总位移
    \item \textbf{判温区}：$T/\Theta_E\gg1$ 用经典近似，$T/\Theta_E\ll1$ 保留零点能
    \item \textbf{验数量级}：最终结果与原子间距、结合能等特征量比较，确认合理
\end{enumerate}
\end{thinkbox}
